\documentclass[11pt]{article}
\usepackage{reusv}
\usepackage{listings}
\usepackage{enumitem}

\lstset{basicstyle=\ttfamily\small, frame=single, breaklines=true}
\setborderthickness{0.5pt}
\setbordermargin{1cm}
\setbordercolor{black}

\slimtitle{A1. 좌표계 및 궤도요소 변환}

\begin{document}
\drawborder
\thispagestyle{firstpage}

\lightertitle{A1. 좌표계 및 궤도요소 변환}
\def\lighterdate{February 12, 2026}
\lighterauthor{Suwon Lee}
\lighteraddress{}{Future Mobility Control Lab, Department of Future Mobility, Kookmin University, Seoul, Republic of Korea}
\lighteremail{suwon@kookmin.ac.kr}

\begin{abstract}
본 문서는 지구 중심 관성 좌표계(ECI)와 고전 궤도요소(classical orbital elements) 사이의 상호 변환에 관한 수학적 배경 및 구현 세부사항을 기술한다.
순변환(OE $\to$ ECI), 역변환(ECI $\to$ OE), CasADi 심볼릭 변환의 세 가지 변환을 다루며,
이들은 LEO-to-LEO 연속 추력 궤도전이 최적화 문제의 경계 조건 설정 및 궤적 후처리에 핵심적으로 사용된다.
\end{abstract}

% ============================================================================
\section{목적}
% ============================================================================

본 문서는 지구 중심 관성 좌표계(ECI)와 고전 궤도요소(classical orbital elements) 사이의 상호 변환에 관한 수학적 배경 및 구현 세부사항을 기술한다. 구체적으로 다음 세 가지 변환을 다룬다.

\begin{enumerate}
  \item \textbf{순변환 (OE $\to$ ECI)}: 6개의 고전 궤도요소로부터 ECI 위치/속도 벡터를 계산한다.
  \item \textbf{역변환 (ECI $\to$ OE)}: ECI 위치/속도 벡터로부터 6개의 궤도요소를 복원한다.
  \item \textbf{CasADi 심볼릭 변환}: 순변환의 CasADi 자동미분 호환 버전을 제공한다.
\end{enumerate}

이 변환들은 LEO-to-LEO 연속 추력 궤도전이 최적화 문제의 경계 조건 설정 및 궤적 후처리에 핵심적으로 사용된다.

대응 코드: \texttt{src/orbit\_transfer/astrodynamics/orbital\_elements.py}

% ============================================================================
\section{수학적 배경}
% ============================================================================

\subsection{ECI 좌표계 정의}

지구 중심 관성 좌표계(Earth-Centered Inertial, ECI)는 지구 중심을 원점으로 하는 관성 기준계이다~\cite{Vallado2013}. 세 축은 다음과 같이 정의된다.

\begin{itemize}
  \item \textbf{$\hat{X}$축}: 춘분점(vernal equinox) 방향. 태양이 천구 적도를 남에서 북으로 통과하는 방향이다.
  \item \textbf{$\hat{Y}$축}: 적도면(equatorial plane) 내에서 $\hat{X}$축에 직교하며, 오른손 좌표계를 형성하는 방향이다. 즉 $\hat{Y} = \hat{Z} \times \hat{X}$이다.
  \item \textbf{$\hat{Z}$축}: 지구 자전축(북극) 방향. 적도면에 수직이다.
\end{itemize}

관성 기준계의 조건으로서, ECI 좌표계의 축 방향은 시간에 대해 고정되어 있다(세차, 장동 등은 무시하거나 별도 보정). 따라서 뉴턴 역학의 운동 방정식을 직접 적용할 수 있다.

\subsection{고전 궤도요소}

케플러 궤도를 기술하는 6개의 고전 궤도요소는 다음과 같다~\cite{Vallado2013, Curtis2020}.

\begin{table}[h]
\centering
\begin{tabular}{cccc}
\toprule
기호 & 명칭 & 정의 & 범위 \\
\midrule
$a$ & 장반경 (semi-major axis) & 타원 궤도의 장축 절반 길이 & $a > 0$ \\
$e$ & 이심률 (eccentricity) & 궤도 형상의 타원 정도. $e=0$이면 원, $0<e<1$이면 타원 & $0 \le e < 1$ \\
$i$ & 경사각 (inclination) & 궤도면과 적도면 사이의 각도 & $0 \le i \le \pi$ \\
$\Omega$ & 승교점경도 (RAAN) & 춘분점에서 승교점까지의 적도면 위 각도 & $0 \le \Omega < 2\pi$ \\
$\omega$ & 근점인수 (argument of periapsis) & 승교점에서 근지점까지의 궤도면 위 각도 & $0 \le \omega < 2\pi$ \\
$\nu$ & 진근점이각 (true anomaly) & 근지점에서 위성 위치까지의 궤도면 위 각도 & $0 \le \nu < 2\pi$ \\
\bottomrule
\end{tabular}
\caption{고전 궤도요소의 정의}
\label{tab:oe_definition}
\end{table}

물리적으로, $(a, e)$는 궤도의 \textbf{크기와 형상}, $(i, \Omega)$는 궤도면의 \textbf{공간적 자세}, $\omega$는 궤도면 내 타원의 \textbf{배향}, $\nu$는 궤도 위의 \textbf{위성 위치}를 각각 결정한다.

\subsection{근지점 좌표계 (Perifocal Frame, PQW)}

근지점 좌표계(perifocal frame)는 궤도면에 고정된 좌표계로, 다음 세 축으로 정의된다~\cite{Curtis2020}.

\begin{itemize}
  \item \textbf{$\hat{P}$축}: 궤도 중심에서 근지점(periapsis) 방향
  \item \textbf{$\hat{Q}$축}: 궤도면 내에서 $\hat{P}$에 직교하며, 위성 운동 방향 쪽
  \item \textbf{$\hat{W}$축}: 궤도 각운동량 벡터 방향, $\hat{W} = \hat{P} \times \hat{Q}$
\end{itemize}

이 좌표계에서 위성의 위치와 속도는 진근점이각 $\nu$만의 함수로 간결하게 표현된다.

\subsection{PQW 좌표에서의 위치 및 속도}

\subsubsection{궤도 방정식}

케플러 궤도에서 위성과 중심체 사이의 거리는 다음과 같다.
\[
  r = \frac{p}{1 + e \cos\nu}
\]
여기서 $p$는 반통경(semi-latus rectum)으로, 장반경 및 이심률과 다음 관계를 가진다.
\[
  p = a(1 - e^2)
\]

\subsubsection{PQW 위치 벡터}

근지점이각 $\nu$에서의 위치 벡터를 PQW 좌표로 표현하면 다음과 같다.
\[
  \mathbf{r}_{PQW} = \begin{bmatrix} r\cos\nu \\ r\sin\nu \\ 0 \end{bmatrix}
\]
$\hat{W}$ 성분이 0인 이유는 위성이 궤도면 위에 있기 때문이다.

\subsubsection{PQW 속도 벡터}

속도 벡터의 유도를 위해, 위치 벡터를 시간에 대해 미분한다. 궤도 운동에서 $\dot{\nu} = h / r^2$ (각운동량 보존)이므로, 시간 미분 대신 $\nu$에 대한 미분과 $\dot{\nu}$의 곱으로 표현할 수 있다.
\[
  \mathbf{v}_{PQW} = \frac{d\mathbf{r}_{PQW}}{dt} = \frac{d\mathbf{r}_{PQW}}{d\nu}\dot{\nu}
\]
$\mathbf{r}_{PQW}$를 $\nu$로 미분하면,
\[
  \frac{d\mathbf{r}_{PQW}}{d\nu} = \begin{bmatrix} \dot{r}\cos\nu - r\sin\nu \\ \dot{r}\sin\nu + r\cos\nu \\ 0 \end{bmatrix} \frac{1}{\dot{\nu}}
\]
여기서 궤도 방정식으로부터
\[
  \dot{r} = \frac{pe\sin\nu}{(1+e\cos\nu)^2}\dot{\nu} = \frac{h}{p}e\sin\nu
\]
이고 $h = \sqrt{\mu p}$이므로, 정리하면 다음과 같다.
\[
  \mathbf{v}_{PQW} = \sqrt{\frac{\mu}{p}} \begin{bmatrix} -\sin\nu \\ e + \cos\nu \\ 0 \end{bmatrix}
\]
이 결과는 속도 벡터가 $\nu$와 궤도 파라미터 $(\mu, p, e)$만으로 결정됨을 보여준다.

\subsection{Perifocal $\to$ ECI 회전 변환}

PQW 좌표계에서 ECI 좌표계로의 변환은 3-1-3 오일러 회전(Euler rotation)으로 구성된다~\cite{Vallado2013}. 이 변환은 ECI 좌표계를 세 번의 회전을 통해 PQW 좌표계와 일치시키는 과정의 역으로 이해할 수 있다.

\subsubsection{회전 순서}

ECI $\to$ PQW 변환은 다음 세 단계의 회전으로 수행된다.

\begin{enumerate}
  \item \textbf{$\hat{Z}$축 기준 $\Omega$ 회전}: ECI의 $\hat{X}$축을 승교점 방향(교선, line of nodes)으로 정렬한다.
  \item \textbf{$\hat{X}'$축 기준 $i$ 회전}: 적도면을 궤도면으로 기울인다.
  \item \textbf{$\hat{Z}''$축 기준 $\omega$ 회전}: 교선 방향을 근지점 방향($\hat{P}$)으로 정렬한다.
\end{enumerate}

따라서 ECI $\to$ PQW 변환 행렬은 다음과 같다.
\[
  \mathbf{R}_{ECI \to PQW} = \mathbf{R}_3(\omega) \, \mathbf{R}_1(i) \, \mathbf{R}_3(\Omega)
\]
여기서 기본 회전 행렬들은 다음과 같다.
\[
  \mathbf{R}_3(\theta) = \begin{bmatrix} \cos\theta & \sin\theta & 0 \\ -\sin\theta & \cos\theta & 0 \\ 0 & 0 & 1 \end{bmatrix}, \quad \mathbf{R}_1(\theta) = \begin{bmatrix} 1 & 0 & 0 \\ 0 & \cos\theta & \sin\theta \\ 0 & -\sin\theta & \cos\theta \end{bmatrix}
\]
우리가 필요한 것은 \textbf{PQW $\to$ ECI} 변환이므로, 역변환(전치)을 취한다.
\[
  \mathbf{R}_{PQW \to ECI} = \mathbf{R}_{ECI \to PQW}^T = \mathbf{R}_3^T(\Omega) \, \mathbf{R}_1^T(i) \, \mathbf{R}_3^T(\omega)
\]
이는 $\mathbf{R}_3(-\Omega) \, \mathbf{R}_1(-i) \, \mathbf{R}_3(-\omega)$와 동일하다.

\subsubsection{회전 행렬 원소 전개}

세 행렬의 곱을 전개하면 PQW $\to$ ECI 회전 행렬 $\mathbf{R}$의 원소는 다음과 같다.
\[
  \mathbf{R} = \begin{bmatrix} R_{11} & R_{12} & R_{13} \\ R_{21} & R_{22} & R_{23} \\ R_{31} & R_{32} & R_{33} \end{bmatrix}
\]
각 원소를 명시하면,
\begin{align*}
  R_{11} &= \cos\Omega\cos\omega - \sin\Omega\sin\omega\cos i \\
  R_{12} &= -\cos\Omega\sin\omega - \sin\Omega\cos\omega\cos i \\
  R_{13} &= \sin\Omega\sin i \\
  R_{21} &= \sin\Omega\cos\omega + \cos\Omega\sin\omega\cos i \\
  R_{22} &= -\sin\Omega\sin\omega + \cos\Omega\cos\omega\cos i \\
  R_{23} &= -\cos\Omega\sin i \\
  R_{31} &= \sin\omega\sin i \\
  R_{32} &= \cos\omega\sin i \\
  R_{33} &= \cos i
\end{align*}

\subsubsection{최종 변환}

ECI 좌표에서의 위치/속도 벡터는 다음과 같이 계산된다.
\[
  \mathbf{r}_{ECI} = \mathbf{R} \, \mathbf{r}_{PQW}, \quad \mathbf{v}_{ECI} = \mathbf{R} \, \mathbf{v}_{PQW}
\]
이때 $\mathbf{R}$은 위치/속도에 동일하게 적용된다. 회전은 좌표계만 변환하며, 미분 구조(시간 미분)에 영향을 주지 않기 때문이다.

\subsection{역변환: ECI $\to$ 궤도요소}

ECI 상태벡터 $(\mathbf{r}, \mathbf{v})$로부터 궤도요소를 복원하는 절차는 다음과 같다~\cite{Vallado2013, Curtis2020}.

\subsubsection*{(1) 각운동량 벡터}
\[
  \mathbf{h} = \mathbf{r} \times \mathbf{v}, \quad h = \|\mathbf{h}\|
\]

\subsubsection*{(2) 교선 벡터 (node vector)}
\[
  \mathbf{n} = \hat{K} \times \mathbf{h}, \quad n = \|\mathbf{n}\|
\]
여기서 $\hat{K} = [0, 0, 1]^T$은 ECI의 $\hat{Z}$축 단위벡터이다. $\mathbf{n}$은 승교점 방향을 가리킨다.

\subsubsection*{(3) 이심률 벡터}

이심률 벡터는 근지점 방향을 가리키며, 크기가 이심률과 같다~\cite{Battin1999}.
\[
  \mathbf{e} = \frac{1}{\mu}\left[\left(v^2 - \frac{\mu}{r}\right)\mathbf{r} - (\mathbf{r} \cdot \mathbf{v})\mathbf{v}\right], \quad e = \|\mathbf{e}\|
\]
이 식은 운동 방정식으로부터 유도되는 라플라스-룽게-렌츠(Laplace-Runge-Lenz) 벡터의 무차원 형태이다.

\subsubsection*{(4) 장반경 (vis-viva equation)}

궤도 에너지(specific orbital energy)는 다음과 같다.
\[
  \mathcal{E} = \frac{v^2}{2} - \frac{\mu}{r}
\]
vis-viva 관계로부터 장반경을 구한다.
\[
  a = -\frac{\mu}{2\mathcal{E}}
\]

\subsubsection*{(5) 경사각}
\[
  i = \arccos\left(\frac{h_z}{h}\right)
\]
여기서 $h_z$는 $\mathbf{h}$의 $z$성분이다. $\mathbf{h}$가 $\hat{Z}$축과 이루는 각도가 경사각이다.

\subsubsection*{(6) 승교점경도}
\[
  \Omega = \arccos\left(\frac{n_x}{n}\right)
\]
사분면 판별: $n_y < 0$이면 $\Omega = 2\pi - \Omega$이다. 이는 $\mathbf{n}$이 $\hat{X}$축 아래쪽(음의 $y$방향)에 있으면 $\Omega$가 $\pi$보다 크기 때문이다.

\subsubsection*{(7) 근점인수}
\[
  \omega = \arccos\left(\frac{\mathbf{n} \cdot \mathbf{e}}{n \, e}\right)
\]
사분면 판별: $e_z < 0$이면 $\omega = 2\pi - \omega$이다. 근지점이 궤도면 아래(남쪽)에 있으면 $\omega$는 $\pi$보다 크다.

\subsubsection*{(8) 진근점이각}
\[
  \nu = \arccos\left(\frac{\mathbf{e} \cdot \mathbf{r}}{e \, r}\right)
\]
사분면 판별: $\mathbf{r} \cdot \mathbf{v} < 0$이면 $\nu = 2\pi - \nu$이다. 위성이 원점에서 멀어지는 중($\mathbf{r} \cdot \mathbf{v} > 0$)이면 $\nu < \pi$이고, 가까워지는 중이면 $\nu > \pi$이다.

\subsection{특이값 처리}

고전 궤도요소는 특정 궤도 형상에서 정의가 불분명해지는 특이점(singularity)을 가진다~\cite{Vallado2013}.

\subsubsection{원궤도 ($e = 0$): $\omega$ 미정}

이심률이 0이면 근지점이 존재하지 않으므로 $\omega$가 정의되지 않는다. 이 경우 $\omega = 0$으로 설정하고, $\nu$를 \textbf{argument of latitude} $u$로 대체한다.
\[
  u = \omega + \nu
\]
$\omega = 0$이므로 $u = \nu$가 되며, $\nu$는 승교점에서 위성 위치까지의 궤도면 위 각도를 의미한다.

구체적으로, $\nu$는 다음과 같이 계산한다.
\[
  \nu = \arccos\left(\frac{\hat{\mathbf{n}} \cdot \hat{\mathbf{r}}}{1}\right)
\]
사분면 판별은 $(\mathbf{n} \times \mathbf{r}) \cdot \mathbf{h}$의 부호로 수행한다. 이 외적이 $\mathbf{h}$와 반대 방향이면 $\nu = 2\pi - \nu$이다.

\subsubsection{적도궤도 ($i = 0$): $\Omega$ 미정}

경사각이 0이면 궤도면이 적도면과 일치하여 교선이 존재하지 않으므로 $\Omega$가 정의되지 않는다. 이 경우 $\Omega = 0$으로 설정한다.

$e > 0$인 적도 타원궤도에서는 $\omega$를 이심률 벡터의 방향으로부터 결정한다.
\[
  \omega = \text{atan2}(e_y, e_x)
\]

\subsubsection{원형 적도궤도 ($e = 0$, $i = 0$)}

$\omega$와 $\Omega$가 모두 미정이다. 이 경우 $\Omega = \omega = 0$으로 설정하고, $\nu$를 \textbf{true longitude} $\lambda$로 대체한다.
\[
  \lambda = \Omega + \omega + \nu
\]
$\Omega = \omega = 0$이므로 $\lambda = \nu$이며, 이는 춘분점에서 위성 위치까지의 적도면 위 각도이다.
\[
  \nu = \text{atan2}(r_y, r_x)
\]
이 처리는 관례(convention)에 해당하며, 순변환과 역변환에서 동일한 관례를 사용하면 round-trip 일관성이 보장된다.

\subsection{CasADi 심볼릭 구현}

궤적 최적화 문제에서는 경계 조건으로서 궤도요소-상태벡터 변환이 NLP(Nonlinear Programming) 제약 조건에 포함된다. 이 제약 조건의 기울기(gradient) 및 야코비안(Jacobian)을 효율적으로 계산하기 위해 \textbf{자동미분(automatic differentiation)}이 필요하다.

CasADi는 심볼릭 연산 그래프(symbolic expression graph)를 구축하여 순방향/역방향 자동미분을 지원하는 프레임워크이다~\cite{Andersson2019}. \texttt{oe\_to\_rv\_casadi} 함수는 \texttt{oe\_to\_rv}와 동일한 수학을 CasADi의 \texttt{MX} 심볼릭 타입으로 구현한다.

주요 차이점은 다음과 같다.

\begin{itemize}
  \item \texttt{np.cos/sin} $\to$ \texttt{ca.cos/sin}: CasADi 심볼릭 삼각함수
  \item \texttt{np.array} $\to$ \texttt{ca.vertcat}: 심볼릭 벡터 결합
  \item \texttt{np.sqrt} $\to$ \texttt{ca.sqrt}: 심볼릭 제곱근
  \item 행렬 구성: \texttt{ca.horzcat}과 \texttt{ca.vertcat}의 조합으로 $3 \times 3$ 회전 행렬 생성
\end{itemize}

이렇게 구성된 심볼릭 표현은 \texttt{ca.Function}으로 래핑한 후 IPOPT 등의 NLP 솔버에 직접 전달할 수 있다. CasADi가 내부적으로 연쇄 법칙(chain rule)을 적용하여 야코비안 $\partial \mathbf{r}/\partial \mathbf{oe}$, $\partial \mathbf{v}/\partial \mathbf{oe}$를 정확하게 계산한다.

역변환 \texttt{rv\_to\_oe}는 \texttt{arccos}, \texttt{clip}, 조건 분기 등 CasADi에서 직접 미분 불가능한 연산을 포함하므로, 심볼릭 버전을 제공하지 않는다. 최적화 문제에서는 순변환만 제약 조건에 사용되기 때문에 이것으로 충분하다.

% ============================================================================
\section{구현 매핑}
% ============================================================================

\subsection{함수 매핑 테이블}

\begin{table}[h]
\centering
\begin{tabular}{lllll}
\toprule
수학 표현 & Python 함수 & 파일 위치 & 백엔드 & 비고 \\
\midrule
OE $\to$ $(\mathbf{r}, \mathbf{v})$ & \texttt{oe\_to\_rv} & \texttt{orbital\_elements.py:7-70} & NumPy & 수치 평가용 \\
$(\mathbf{r}, \mathbf{v})$ $\to$ OE & \texttt{rv\_to\_oe} & \texttt{orbital\_elements.py:73-170} & NumPy & 특이값 분기 포함 \\
OE $\to$ $(\mathbf{r}, \mathbf{v})$ & \texttt{oe\_to\_rv\_casadi} & \texttt{orbital\_elements.py:173-238} & CasADi & NLP 자동미분용 \\
\bottomrule
\end{tabular}
\caption{함수 매핑 테이블}
\label{tab:function_mapping}
\end{table}

\subsection{순변환 (\texttt{oe\_to\_rv}) 구현 상세}

\begin{table}[h]
\centering
\begin{tabular}{lll}
\toprule
수학 & 코드 & 행 \\
\midrule
$p = a(1-e^2)$ & \texttt{p = a * (1.0 - e**2)} & 33 \\
$r = p/(1+e\cos\nu)$ & \texttt{r\_mag = p / (1.0 + e * np.cos(nu))} & 36 \\
$\mathbf{r}_{PQW}$ & \texttt{r\_pqw = np.array([r\_mag*np.cos(nu), r\_mag*np.sin(nu), 0.0])} & 39--41 \\
$\mathbf{v}_{PQW}$ & \texttt{v\_pqw = np.sqrt(mu/p) * np.array([-np.sin(nu), e+np.cos(nu), 0.0])} & 43--45 \\
$\mathbf{R}$ 행렬 원소 & \texttt{R = np.array([[...], [...], [...]])} & 55--65 \\
$\mathbf{r} = \mathbf{R}\mathbf{r}_{PQW}$ & \texttt{r = R @ r\_pqw} & 67 \\
$\mathbf{v} = \mathbf{R}\mathbf{v}_{PQW}$ & \texttt{v = R @ v\_pqw} & 68 \\
\bottomrule
\end{tabular}
\caption{순변환 \texttt{oe\_to\_rv} 구현 상세}
\label{tab:forward_impl}
\end{table}

\subsection{역변환 (\texttt{rv\_to\_oe}) 구현 상세}

\begin{table}[h]
\centering
\begin{tabular}{lll}
\toprule
수학 & 코드 & 행 \\
\midrule
$\mathbf{h} = \mathbf{r} \times \mathbf{v}$ & \texttt{h = np.cross(r, v)} & 99 \\
$\mathbf{n} = \hat{K} \times \mathbf{h}$ & \texttt{n = np.cross(K, h)} & 104 \\
이심률 벡터 & \texttt{e\_vec = ((v\_mag**2 - mu/r\_mag)*r - np.dot(r,v)*v) / mu} & 108 \\
$\mathcal{E} = v^2/2 - \mu/r$ & \texttt{energy = 0.5*v\_mag**2 - mu/r\_mag} & 112 \\
$a = -\mu/(2\mathcal{E})$ & \texttt{a = -mu / (2.0*energy)} & 113 \\
$i = \arccos(h_z/h)$ & \texttt{i = np.arccos(np.clip(h[2]/h\_mag, -1, 1))} & 116 \\
$e=0$ 분기 & \texttt{if e < 1e-10:} & 131 \\
$i=0$ 분기 & \texttt{if i < 1e-10:} & 119, 135 \\
\bottomrule
\end{tabular}
\caption{역변환 \texttt{rv\_to\_oe} 구현 상세}
\label{tab:inverse_impl}
\end{table}

\subsection{특이값 처리 분기 구조}

\begin{lstlisting}
rv_to_oe
+-- i < 1e-10?  -> Omega = 0
|   +-- else    -> Omega = arccos(n_x/n), quadrant correction
+-- e < 1e-10?  -> omega = 0
|   +-- i < 1e-10?  -> nu = atan2(r_y, r_x)     [true longitude]
|   +-- else        -> nu = arccos(n_hat . r_hat) [argument of latitude]
+-- else (e > 0)
    +-- i < 1e-10?  -> omega = atan2(e_y, e_x)   [longitude of periapsis]
    +-- else        -> omega = arccos(n.e / ne), quadrant correction
                       nu = arccos(e.r / er), quadrant correction
\end{lstlisting}

% ============================================================================
\section{수치 검증}
% ============================================================================

수치 검증은 \texttt{tests/test\_orbital\_elements.py}에 구현되어 있다. 검증 항목은 다음 네 가지이다.

\subsection{Round-trip 일관성 (OE $\to$ RV $\to$ OE)}

다양한 궤도 유형에 대해 순변환 후 역변환을 수행하고, 원래 궤도요소와의 차이를 확인한다.

\begin{table}[h]
\centering
\begin{tabular}{lcccc}
\toprule
테스트 케이스 & $a$ [km] & $e$ & $i$ [deg] & 허용 오차 \\
\midrule
LEO-ISS & $R_E + 400$ & 0.001 & 51.6 & $\Delta a < 10^{-8}$, $\Delta e < 10^{-12}$ \\
High eccentricity & $R_E + 500$ & 0.3 & 28.5 & 각도 차이 $< 10^{-12}$ rad \\
GEO-like & 42164 & 0.01 & 0.05 & 동일 \\
Molniya-like & 26600 & 0.74 & 63.4 & 동일 \\
\bottomrule
\end{tabular}
\caption{Round-trip 일관성 검증 테스트 케이스}
\label{tab:roundtrip_cases}
\end{table}

각도 비교 시 $2\pi$ 주기성을 고려한 차이 함수를 사용한다.
\[
  \Delta\theta = \min\left(|\theta_1 - \theta_2| \bmod 2\pi, \; 2\pi - |\theta_1 - \theta_2| \bmod 2\pi\right)
\]

\subsection{원궤도 특이 케이스 ($e = 0$)}

\begin{itemize}
  \item \textbf{원형 경사궤도} ($e=0$, $i=51.6^\circ$): 궤도 반경 $\|\mathbf{r}\| = a$, 원궤도 속력 $\|\mathbf{v}\| = \sqrt{\mu/a}$ 확인. 역변환 후 $e < 10^{-10}$ 확인.
  \item \textbf{원형 적도궤도} ($e=0$, $i=0$): 동일한 반경/속력 검증. 역변환 후 $e < 10^{-10}$, $i < 10^{-10}$ 확인.
\end{itemize}

\subsection{적도궤도 특이 케이스 ($i = 0$)}

\begin{itemize}
  \item \textbf{적도 타원궤도} ($e=0.1$, $i=0$): ECI $z$성분이 0임을 확인 ($|r_z| < 10^{-10}$, $|v_z| < 10^{-10}$). 역변환 후 $\omega + \nu$의 합이 보존됨을 확인.
\end{itemize}

\subsection{CasADi/NumPy 일관성}

동일한 입력에 대해 \texttt{oe\_to\_rv}(NumPy)와 \texttt{oe\_to\_rv\_casadi}(CasADi)의 출력 차이를 비교한다.

\begin{table}[h]
\centering
\begin{tabular}{ll}
\toprule
테스트 케이스 & 허용 오차 \\
\midrule
LEO-ISS & $\max|\mathbf{r}_{CA} - \mathbf{r}_{NP}| < 10^{-14}$ \\
High eccentricity & $\max|\mathbf{v}_{CA} - \mathbf{v}_{NP}| < 10^{-14}$ \\
GEO circular ($e=0$) & 동일 \\
\bottomrule
\end{tabular}
\caption{CasADi/NumPy 일관성 검증}
\label{tab:casadi_consistency}
\end{table}

$10^{-14}$ 수준의 일치는 두 구현이 동일한 수학 연산을 수행하며, 부동소수점 연산 순서까지 일관적임을 의미한다.

% ============================================================================
\section{참고문헌}
% ============================================================================

\bibliographystyle{IEEEtran}
\bibliography{references}

\end{document}
